\chapter{Matriz de consistencia}

\begin{table}[H]
\centering
\caption{Matriz de consistencia: problema, objetivos, hip\'otesis, variables y m\'etodo}
\label{tab:anexo_consistencia}
\footnotesize
\begin{tabularx}{\textwidth}{p{2.2cm}X}
\toprule
\textbf{Elemento} & \textbf{Descripci\'on} \\
\midrule
Problema & Los productores de palta Hass de La Libertad toman decisiones de siembra sin informaci\'on anticipada sobre precios FOB, lo que incrementa el riesgo econ\'omico ante la volatilidad del mercado exportador. \\
\midrule
Objetivo general & Desarrollar un modelo predictivo de precios FOB de palta Hass que apoye la toma de decisiones de siembra en productores de La Libertad. \\
\midrule
Objetivos espec\'ificos & (1) Caracterizar la evoluci\'on del precio FOB; (2) Identificar variables explicativas; (3) Comparar SARIMA, LSTM y SVR; (4) Seleccionar el mejor modelo; (5) Desplegar aplicaci\'on m\'ovil con simulador de rentabilidad. \\
\midrule
Hip\'otesis general (Hg) & El modelo predictivo impacta positivamente en el apoyo a la decisi\'on de siembra al mejorar pron\'ostico, rentabilidad, escenarios, an\'alisis hist\'orico y reportes. \\
\midrule
Hip\'otesis espec\'ificas & He1 (rentabilidad), He2 (escenarios), He3 (precisi\'on predictiva), He4 (an\'alisis hist\'orico), He5 (reportes de ciclos). \\
\midrule
Variables & Precio FOB (dependiente); volumen exportado, estacionalidad, rezagos y destino (independientes); utilidad percibida y m\'etricas de error (evaluaci\'on). \\
\midrule
M\'etodo & Dise\~no mixto: cuasiexperimental de series temporales (SARIMA, LSTM, SVR) y preexperimental con pretest--postest en productores (CRISP-DM). \\
\midrule
T\'ecnicas & An\'alisis documental, ETL en Python, modelado predictivo, evaluaci\'on con RMSE/MAPE/$R^{2}$, cuestionario de utilidad percibida. \\
\bottomrule
\end{tabularx}
\parbox{0.94\textwidth}{\small \textit{Nota.} Elaboraci\'on propia con base en los Cap\'itulos I y II.}
\end{table}
